% MultiSpec-GeoDiff — Project Proposal Submission
% Compile: latexmk -xelatex submission.tex
\documentclass[10pt,a4paper]{ctexart}
\usepackage{amsmath,amssymb,bm}
\usepackage{graphicx}
\usepackage{booktabs,tabularx}
\usepackage[margin=1.2cm]{geometry}
\usepackage{hyperref}
\usepackage{caption}
\usepackage{enumitem}
\usepackage{xcolor}

\hypersetup{
  colorlinks=true,
  linkcolor=black,
  citecolor=black,
  urlcolor=black,
  pdftitle={MultiSpec-GeoDiff：多模态谱图驱动的分子结构反演},
}

\setlength{\parindent}{0pt}
\setlength{\parskip}{1.5pt plus 0.5pt}
\linespread{0.92}
\captionsetup{font=small,labelfont=small,skip=1pt}

\setlength{\abovedisplayskip}{3pt}
\setlength{\belowdisplayskip}{3pt}
\setlength{\abovedisplayshortskip}{2pt}
\setlength{\belowdisplayshortskip}{2pt}

\newcommand{\spec}{\mathcal S}
\newcommand{\R}{\mathbb R}


\begin{document}

\thispagestyle{empty}

{\LARGE\bfseries MultiSpec-GeoDiff：多模态谱图驱动的分子结构反演\par}
\vspace{3pt}
{\small\bfseries 谢翔宇\quad 中国科学技术大学\quad \href{mailto:xxy220348@mail.ustc.edu.cn}{xxy220348@mail.ustc.edu.cn}\par}
\vspace{4pt}
\hrule height 0.6pt
\vspace{4pt}

% ═══════════════════════════════════════════════════════════════
% NOTICE: Document structure explanation
% ═══════════════════════════════════════════════════════════════
\vspace{2pt}
\noindent\fbox{%
\begin{minipage}{\dimexpr\textwidth-2\fboxsep-2\fboxrule\relax}
\small
\textbf{说明：}第 1 节为 500 字内项目提议正文，严格对应``具体目标---核心价值---可行性简析''三项要求；第 2 节及后续数学结构、架构图与阶段路线为补充技术说明，用于解释 Demo 与长期研究路线的关系，不计入项目提议正文字数。\par
\textbf{当前提交 Demo 实现的是 IR-only Stage-I 可行性闭环}；MultiSpec 与 GeoDiff 表示完整研究方向，多模态融合、图扩散、pairwise-distance 预测与 TFN/SE(3) 精修为后续扩展，不作为当前已完成能力宣称。
\end{minipage}%
}
\vspace{6pt}

% ═══════════════════════════════════════════════════════════════
% SECTION 1: OFFICIAL PROJECT PROPOSAL (≤500 chars)
% ═══════════════════════════════════════════════════════════════
\section*{一、项目提议正文（500 字内）}

\textbf{具体目标：}
拟构建多模态谱图驱动的几何分子结构反演模型，将反演形式化为几何约束下的多模态条件生成问题：
\[
\spec_{\mathrm{multi}} \rightarrow (c,C) \rightarrow G_{2D} \rightarrow D \rightarrow R^{(0)} \rightarrow (R,\mathcal T) \rightarrow \hat{\spec} \rightarrow \mathrm{TopK}.
\]
当前 Demo 使用 200 条 NIST 实验 IR 光谱，验证 IR 谱图读取、坐标感知编码、候选召回、前向谱图重排与 Top-K 可视化闭环；完整阶段可扩展至 QM9S 与公开多模态谱图数据库。

\textbf{核心价值：}
谱图反演是一谱多解的逆问题---横坐标携带物理含义，分子具有图拓扑与三维几何。本方案以横坐标偏置注意力尊重物理坐标，以可增删节点图扩散自适应分子大小，以 pairwise distance 矩阵为 2D--3D 桥梁，以 TFN-Transformer 多阶张量特征与张量内积注意力实现 SE(3) 等变建模，以前向谱图一致性验证候选，将反演从序列翻译提升至生成与验证闭环。

\textbf{可行性简析：}
分三阶段推进---阶段一验证 IR 编码与候选重排闭环（已实现）；阶段二引入图扩散与距离预测；阶段三加入 parity-aware TFN-Transformer 与前向谱图闭环。当前 Demo 使用 200 条 NIST 实验 IR 光谱，验证 IR 编码---候选召回---前向谱图重排---Top-K 输出闭环。主要风险为实验噪声、一谱多解与数据规模有限；验证方式包括 Top-K 命中率、谱图相似度、候选重排增益、分子式/官能团匹配率、trace log 与 CI 测试。

\vspace{2pt}
\noindent\rule{\textwidth}{0.4pt}
\vspace{1pt}
{\small\noindent\textbf{（以上为项目提议正文，以下为补充技术说明，不计入 500 字限制。）}}
\vspace{6pt}

% ═══════════════════════════════════════════════════════════════
% SUPPLEMENTARY FIGURE
% ═══════════════════════════════════════════════════════════════
{\centering
\includegraphics[width=\textwidth,height=0.42\textheight,keepaspectratio]{figures/fig0.png}
\captionof{figure}{图 1（补充说明）.\quad MultiSpec-GeoDiff 总体架构。实线模块表示当前 Stage-I 已实现的 IR 检索/重排闭环（200 条 NIST 实验 IR 光谱）；虚线模块表示 Stage-II/III 的图扩散、距离矩阵和 TFN-Transformer 几何扩展路线。}\par}

\vspace{4pt}
\hrule height 0.4pt
\vspace{4pt}

% ═══════════════════════════════════════════════════════════════
% SECTION 2: SUPPLEMENTARY — CORE MATH FORMULATION
% ═══════════════════════════════════════════════════════════════
\section*{补充技术说明：核心数学结构}

\textbf{输入表示。}
多模态谱图表为峰集合 $S^{(m)} = \{(\nu_i^{(m)}, I_i^{(m)})\}_{i=1}^{L_m}$，其中 $m \in \{\mathrm{IR}, \mathrm{Raman}, \mathrm{NMR}, \mathrm{MS}, \mathrm{UV}, \ldots\}$。当前 Stage-I 仅使用 IR 单模态；Raman、NMR、MS、UV、CD/VCD/ROA 为未来扩展模态。

\vspace{2pt}
\textbf{1. 坐标感知谱图编码（Stage-I 已实现）}
\begin{equation}\label{eq:embed}
H_i^{(m,0)} = \mathrm{CNN}_m(S_i^{(m)}) + e_m + p_m(\nu_i^{(m)}),\qquad
p_m(\nu) = \mathrm{MLP}_m[\nu,\;\sin(\omega_1\nu),\;\cos(\omega_1\nu),\;\ldots].
\end{equation}
模态内自注意力引入真实横坐标距离偏置，使模型感知峰间物理间距：
\begin{equation}\label{eq:attn}
A_{ij}^{(m)} = \mathrm{softmax}_j\!\left(\frac{Q_i^{(m)} K_j^{(m)\top}}{\sqrt d} + b_m(|\nu_i^{(m)} - \nu_j^{(m)}|)\right).
\end{equation}
跨模态注意力不引入横坐标偏置（不同模态坐标系不可公度），仅学习证据兼容性：
\[
A_{ij}^{(m\leftarrow n)} = \mathrm{softmax}_j\!\left(\frac{Q_i^{(m)} K_j^{(n)\top}}{\sqrt d} + \gamma_{mn}\right).
\]
编码器输出全局条件 $c$ 与细粒度 token 条件 $C$。\textbf{式 (\ref{eq:embed})--(\ref{eq:attn}) 为当前 Stage-I 已实现模块；跨模态注意力为未来多模态扩展所用。}

\vspace{2pt}
\textbf{2. 二维图扩散（Stage-II roadmap）}
\[
G_t = (X_t, E_t),\quad a_t \in \{\mathrm{insert}, \mathrm{delete}, \mathrm{atom}, \mathrm{bond}\},\quad
D_\theta(G_t, t, c, C) \rightarrow a_t.
\]
反扩散 Graph Transformer 引入拓扑距离偏置与键类型偏置：
\[
\alpha_{ij} = \mathrm{softmax}_j\!\left(\frac{Q_i K_j^\top}{\sqrt d} + \beta_{\mathrm{SPD}(i,j)} + \eta_{e_{ij}}\right).
\]
输出 $G_{2D} = (X, E)$---二维图给出原子与化学键，但光谱响应依赖三维几何。\textbf{图扩散为 Stage-II roadmap，当前 Demo 未包含训练完成的图扩散模型。}

\vspace{2pt}
\textbf{3. Pairwise distance：2D 到 3D 的桥梁（Stage-II roadmap）}
\[
\hat D = f_\theta(G_{2D}, c, C),\qquad
R^{(0)} = \mathrm{DistGeom}(\hat D),\qquad
\mathcal L_D = \sum\nolimits_{i<j} |\hat D_{ij} - D_{ij}^{\mathrm{ref}}|^2.
\]
绝对坐标具有旋转平移规范歧义，pairwise 距离矩阵天然不变：$D_{ij} = \|r_i - r_j\|$。\textbf{距离矩阵预测为 Stage-II roadmap，当前 Demo 未包含训练完成的 pairwise 距离预测器。}

\vspace{2pt}
\textbf{4. TFN-Transformer：坐标与张量等变精修（Stage-III roadmap）}
\begin{equation}\label{eq:tfn_attn}
s_{ij} = \sum_{l,p,c} \langle Q_{i,c}^{(l,p)}, K_{j,c}^{(l,p)} \rangle
        + b_{\mathrm{dist}}(\|r_i - r_j\|) + b_{\mathrm{edge}}(e_{ij}),\qquad
\alpha_{ij} = \mathrm{softmax}_j(s_{ij}).
\end{equation}
TFN 消息通过球谐张量积实现等变特征交互：
\[
M_{ij}^{(l_3,p_3)} = \sum_{l_1,p_1,l_2,p_2} \phi(\|r_{ij}\|, e_{ij})
                     \bigl(T_j^{(l_1,p_1)} \otimes Y_{l_2}^{(p_2)}(\hat r_{ij})\bigr)_{l_3,p_3},\quad
p_3 = p_1 p_2.
\]
坐标以 EGNN 风格相对更新：
\begin{equation}\label{eq:coord}
r_i^{(\ell+1)} = r_i^{(\ell)} + \sum_{j \neq i} \alpha_{ij} (r_i^{(\ell)} - r_j^{(\ell)})\,
                 \phi_x(\|r_i - r_j\|^2, e_{ij}, \langle\mathcal T_i, \mathcal T_j\rangle_{\mathrm{tensor}}, c).
\end{equation}
\textbf{式 (\ref{eq:tfn_attn})--(\ref{eq:coord}) 为 Stage-III roadmap，当前 Demo 未包含训练完成的 TFN-Transformer 模型。}

\vspace{2pt}
\textbf{5. 手性与宇称（Stage-III roadmap）}
三重向量积 $\chi_{ijkl} = (r_j - r_i) \cdot [(r_k - r_i) \times (r_l - r_i)]$ 在旋转下不变、反射下反号，属奇宇称标量（$0o$）通道，可在 CD/VCD/ROA 等手性敏感模态监督下编码手性信息。当前 Demo 未实现手性区分。

\vspace{2pt}
\textbf{6. 前向谱图一致性与 Top-K 后验重排（Stage-I 已实现重排逻辑）}
\[
\hat{\spec}_k = F_{\mathrm{spec}}(G_{2D,k}, R_k, \mathcal T_k),\qquad
s_k = -\mathrm{Dist}(\spec, \hat{\spec}_k) + \lambda_{\mathrm{chem}} S_{\mathrm{chem}}(G_k) + \lambda_{\mathrm{traj}} S_{\mathrm{diff}}(G_k).
\]
最终 $\mathrm{TopK} = \mathrm{argsort}_k(s_k)$。当前 Stage-I 使用检索候选 + 前向谱图三指标重排（余弦相似度 + L1 距离 + 峰匹配）；未来阶段可替换为谱图条件图扩散生成与前向 TFN 谱图预测。生成模块负责提出候选，前向谱图模块负责物理校验，二者共同形成谱图反演的闭环。

\vspace{4pt}
\hrule height 0.4pt
\vspace{4pt}

% ═══════════════════════════════════════════════════════════════
% SECTION 3: SUPPLEMENTARY — DEMO ALIGNMENT & STAGE ROADMAP
% ═══════════════════════════════════════════════════════════════
\section*{补充说明：Demo 对齐与阶段路线}

\textbf{Stage-I（已实现）：}
当前 Demo 使用 200 条 NIST 实验 IR 光谱，实现 IR 谱图预处理、坐标感知编码、候选检索、前向谱图重排与 Top-K 可视化输出，形成可追溯、可复现的最小闭环。22 个 pytest 测试全部通过，输出含完整 trace log（\texttt{05\_outputs/trace\_log.json}），验证方式包括 Top-K 命中率、谱图相似度、候选重排增益、分子式/官能团匹配率与 CI 测试。

\textbf{Stage-II（近期 roadmap）：}
引入大小自适应图扩散（insert/delete）生成二维分子图候选；pairwise 距离矩阵预测头联合训练，经 distance geometry 恢复初始三维坐标。

\textbf{Stage-III（远期）：}
引入 parity-aware TFN-Transformer 等变精修与多阶张量前向谱图预测；扩展至 CD/VCD/ROA 等手性敏感模态与 Raman、NMR、MS、UV 等多模态谱图。

\vspace{4pt}
\hrule height 0.4pt
\vspace{3pt}
{\small
受真实多模态实验谱图数据与量化标注成本限制，当前 Demo 以 200 条 NIST 实验 IR 光谱验证可执行闭环，其目标是证明技术路径可行性而非报告生产级性能。完整研究计划需依赖更大规模公开/实验谱图数据库与量化弱监督。\\
项目仓库：\url{https://github.com/techandscixie2005/MultiSpec-GeoDiff-Demo}\quad
Demo 说明：\texttt{02\_demo\_document/DEMO说明.md}、\texttt{02\_demo\_document/MultiSpec-GeoDiff\_Demo\_说明.pdf}.
}

\end{document}
